% ===== 原文 pp.91–104（PDF p.106–119）· 第 3 章 §3.1–3.2 · tr-06 =====
% [原文 p.91]
\section{对数几何}

对数几何由 Illusie 与 Fontaine [59] 以及 K.~Kato [65] 引入。对数几何的起源，尤其是“对数”这一术语，来自对数微分。

例如，假设人们想研究一个开代数簇，比如一个非异拟射影簇 $X$。如果 $X$ 含于一个射影簇 $\overline{X}$ 之中，且 $D=\overline{X}\setminus X$ 是正规交叉除子\footnote{正规交叉除子 $D$ 是指这样的除子：在 $\overline{X}$ 任意一点的邻域内存在坐标 $x_1,\ldots,x_n$，使得 $D$ 由 $x_1\cdots x_p=0$ 给出，其中 $p\le n$。}，则 $\Omega^q_{\overline{X}}(\log D)$ 被定义为 $i_*\Omega^q_X$ 的子层（其中 $i:X\hookrightarrow\overline{X}$ 是包含映射），它在局部由
\[
  \frac{\dd x_1}{x_1},\ldots,\frac{\dd x_p}{x_p},\dd x_{p+1},\ldots,\dd x_n
\]
生成，只要 $D$ 在局部由 $x_1\cdots x_p=0$ 给出。当然 $\frac{\dd x_i}{x_i}=\dd\log(x_i)$，于是有了“对数”一词。

这些层最初的重要性由 Deligne 对 $X$ 上混合 Hodge 结构的构造所体现。特别地，外微分把对数形式映为对数形式，而以它为微分的复形 $\Omega^{\bullet}_{\overline{X}}(\log D)$ 实际上计算了 $X$ 的上同调：我们有
\[
  \mathbb{H}^q\bigl(\overline{X},\Omega^{\bullet}_{\overline{X}}(\log D)\bigr)\cong H^q(X,\C).
\]
如同第 2 章，$\mathbb{H}^q$ 表示超上同调。这是给 $H^q(X,\C)$ 赋予混合 Hodge 结构的重要一步。

% [Tutor Note: p.91 核对点：定义里 $\Omega^q_{\overline X}(\log D)$ 是 $i_*\Omega^q_X$ 的“子层”而非 $i_*\Omega^q_{\overline X}$——$\dd x_i/x_i$ 在 $X$ 上有极点，只在 $\overline X$ 上作为亚纯形式有意义；正因如此才能把非紧簇 $X$ 的 de Rham 复形与紧化上的对数复形对上。]
% [Tutor Note: p.91 关于上式的特例：当 $D$ 不是正规交叉（例如有尖点）时，$\mathbb{H}^q\cong H^q(X,\C)$ 一般不成立；原文此处所谓“computes the cohomology”依赖 $D$ 正规交叉的假设，属 Deligne 混合 Hodge 理论的标准结论。]

对数微分层另一个重要应用来自 Steenbrink 对正规交叉退化的极限混合 Hodge 结构的构造。考虑平坦族 $f:\cX\to S$，其中 $S$ 是非异一维概形，且对 $s\in S$，除闭点 $0\in S$ 外 $X_s$ 都是非异的。此外，$f$ 在适当的局部坐标卡中总可写成 $(x_1,\ldots,x_n)\mapsto x_1\cdots x_p$，其中 $p\le n$。于是有相对对数 $q$-形式的层
\[
  \Omega^q_{\cX/S}(\log \cX_0):=\Omega^q_{\cX}(\log \cX_0)/(f^*\Omega^1_S(\log 0)\wedge\Omega^{q-1}_{\cX}(\log \cX_0)).
\]
特别地，$\Omega^1_{\cX/S}(\log \cX_0)$ 在局部由
\[
  \dd\log x_1,\ldots,\dd\log x_p,\dd x_{p+1},\ldots,\dd x_n
\]
生成，模去关系 $\dd\log x_1+\cdots+\dd\log x_p=0$。把 $\Omega^q_{\cX/S}(\log \cX_0)$ 限制到 $\cX_0$ 上得到的层，我们暂时记作 $\Omega^q_{\cX_0^\dagger}$，以区别于通常的微分层 $\Omega^q_{\cX_0}$。现在容易验证 $\Omega^q_{\cX/S}(\log \cX_0)$ 局部自由，因而 $\Omega^q_{\cX_0^\dagger}$ 也局部自由，而 $\Omega^q_{\cX_0}$ 不局部自由。外微分

% [原文 p.92]
仍然有意义，无论在 $\Omega^q_{\cX/S}(\log \cX_0)$ 上还是在 $\Omega^q_{\cX_0^\dagger}$ 上。若 $f$ 紧合，则可以证明（见 [107]）$R^pf_*\Omega^q_{\cX/S}(\log \cX_0)$ 局部自由。在 $\cX_0$ 之外，$\Omega^q_{\cX/S}(\log \cX_0)$ 就是通常的 $q$-形式层，因此在 $0$ 之外，$R^pf_*\Omega^q_{\cX/S}(\log \cX_0)$ 是以 $s$ 处纤维为 Dolbeault 上同调群 $H^p(X_s,\Omega^q_{X_s})$ 的向量丛。$0$ 处的纤维与 $H^p(\cX_0,\Omega^q_{\cX_0^\dagger})$ 一致。这些群是确定与该退化相联系的极限混合 Hodge 结构的数据的一部分。

这里的要点是：对数微分层比通常的微分层性质好得多，并且以某种方式赋予 $\cX_0$ 光滑簇的某些性质。

对数结构的概念是对这一思想的抽象。本质上，它是一种让某些奇异概形表现得仿佛光滑的方式；正如加藤和也（Kazuya Kato）所描述的，对数结构是一种使奇异簇变光滑的“魔法粉末”（poudre magie）。

我们将在 §3.2 给出对数结构的精确定义并探讨其含义。首先简要回顾环面几何。关于这一主题更详细而愉快的介绍，见 Fulton 的书 [27]。

\subsection{环面几何简论}

\subsubsection{幺半群}

我们将频繁使用幺半群与幺半群环。对我们而言，幺半群 $P$ 是一个集合，带有结合且交换的运算以及单位元。该运算 $P\times P\to P$ 通常（但不总是）写成加法，此时单位元写作 $0\in P$。

给定域 $k$，定义幺半群环
\[
  k[P]:=\bigoplus_{p\in P}kz^p,
\]
其中 $z^p$ 是一个符号，乘法是 $k$-双线性的，并由
\[
  z^p\cdot z^{p'}=z^{p+p'}
\]
确定。

标准例子是自然数的加法幺半群
\[
  \N:=\{0,1,2,\ldots\}.
\]
于是 $k[\N]\cong k[x]$，$k[\N^r]\cong k[x_1,\ldots,x_r]$。

幺半群同态是幺半群之间的映射 $f:P\to Q$，满足 $f(0)=0$ 与 $f(p+p')=f(p)+f(p')$。

称幺半群 $P$ 有限生成，如果存在某个 $r$ 及满同态 $\N^r\to P$。

幺半群 $P$ 的（Grothendieck）群是群
\[
  P^{\mathrm{gp}}:=\{p-p'\mid p,p'\in P\}/\sim,
\]
其中 $p-p'$ 是形式符号，等价关系由 $p-p'\sim q-q'$（当 $p+q'=q+p'$ 时）给出。群运算由 $(p-p')+(q-q')=(p+q)-(p'+q')$ 与 $-(p-p')=p'-p$ 给出。

称幺半群 $P$ 是整的，如果典范映射 $P\to P^{\mathrm{gp}}$ 是单射。

可以用同余关系作幺半群的商：给定幺半群 $P$，等价关系 $E\subseteq P\times P$ 称为同余关系，如果只要 $(q_1,q_2)\in E$、$p\in P$，就有 $(q_1+p,q_2+p)\in E$。同余关系的等价类集合容易看出从 $P$ 继承幺半群结构。若 $Q\subseteq P$ 是子幺半群，

% [原文 p.93]
则 $Q$ 通过 $a\sim b$（若存在 $q,r\in Q$ 使 $a+q=b+r$）在 $P$ 上诱导一个同余关系。于是此时商 $P/Q$ 作为幺半群有意义。

给定幺半群同态 $f_i:Q\to P_i$（$i=1,2$），可以构造纤维余积 $P_1\oplus_Q P_2$，它是 $P_1\oplus P_2$ 按如下同余关系取商：$(p_1,p_2)\sim(p_1',p_2')$，如果存在 $q,r\in Q$ 使得
\[
  p_1+f_1(r)=p_1'+f_1(q)\quad\text{且}\quad p_2+f_2(q)=p_2'+f_2(r).
\]

关于幺半群的许多其他事实，见 [88]。

\subsubsection{来自扇的环面簇}

本节固定数据
\[
  \begin{aligned}
    M&:=\Z^n, & N&:=\operatorname{Hom}_\Z(M,\Z),\\
    M_\R&:=M\otimes_\Z\R, & N_\R&:=N\otimes_\Z\R.
  \end{aligned}
\]
设 $\sigma\subseteq M_\R$ 是如定义 1.7 那样的严格凸有理多面体锥，其对偶锥 $\sigma^{\vee}\subseteq N_\R$ 由
\[
  \sigma^{\vee}=\{n\in N_\R\mid \langle n,m\rangle\ge 0\quad \forall m\in\sigma\}
\]
给出，如例 1.31 (1) 那样。则 $\sigma^{\vee}\cap N$ 是幺半群。这是一个标准事实，称为 Gordan 引理，它由 Carathéodory 定理推出，即该幺半群有限生成。因此 $k[\sigma^{\vee}\cap N]$ 是有限生成代数，而
\[
  X_\sigma:=\operatorname{Spec}k[\sigma^{\vee}\cap N]
\]
是由锥 $\sigma$ 定义的仿射环面簇。

\begin{examples}
\begin{enumerate}[label=(\arabic*),leftmargin=*]
\item 设 $\sigma$ 是第一象限。则 $\sigma^{\vee}$ 也是第一象限，且 $\sigma^{\vee}\cap N$ 是幺半群 $\N^2$。此时幺半群环 $k[\sigma^{\vee}\cap N]$ 是多项式环 $k[x,y]$，故 $X_\sigma\cong\mathbb{A}^2_k$。
\item 设 $\sigma$ 是 $\R^2$ 中由 $(1,0)$ 与 $(1,e)$（$e$ 为正整数）生成的锥。则 $\sigma^{\vee}$ 由 $(0,1)$ 与 $(e,-1)$ 生成。但注意幺半群 $\sigma^{\vee}\cap\Z^2$ 并不由 $(0,1)$ 与 $(e,-1)$ 生成，因为 $(1,0)$ 不能表示为这两个向量的整线性组合。事实上，作为幺半群它由 $(0,1)$、$(e,-1)$ 与 $(1,0)$ 生成。这些生成元满足显然的关系 $(0,1)+(e,-1)=e(1,0)$。于是幺半群环同构于 $k[x,y,t]/(xy-t^e)$，其中 $x=z^{(0,1)}$，$y=z^{(e,-1)}$，$t=z^{(1,0)}$。所以 $V_\sigma$ 是带有 $A_{e-1}$ 奇点的曲面。
\end{enumerate}
\end{examples}

% [Tutor Note: p.93 (2) 末句的 $V_\sigma$ 应为 $X_\sigma$：本页上文定义的是 $X_\sigma:=\operatorname{Spec}k[\sigma^\vee\cap N]$，且原文别处一律用 $X_\sigma$；此处 $V_\sigma$ 为原文笔误，照译以存真。]

现在设 $\Sigma$ 是 $M_\R$ 中的扇（如定义 1.7）。则对每个锥 $\sigma\in\Sigma$ 得到仿射环面簇。此外，若 $\tau,\sigma\in\Sigma$ 且 $\tau$ 是 $\sigma$ 的面，则 $\sigma^{\vee}\subseteq\tau^{\vee}$，给出幺半群环的包含。这实际上把 $X_\tau$ 描述为 $X_\sigma$ 的开子集。对 $\sigma_1,\sigma_2\in\Sigma$ 且 $\tau=\sigma_1\cap\sigma_2$，可沿公共开子集 $X_\tau\subseteq X_{\sigma_1}$ 与 $X_\tau\subseteq X_{\sigma_2}$ 粘接 $X_{\sigma_1}$ 与 $X_{\sigma_2}$。完成这些粘接后，得到环面簇 $X_\Sigma$，它是 $k$ 上的分离代数簇。

扇 $\Sigma$ 的性质反映在簇 $X_\Sigma$ 的性质中。例如，
\begin{enumerate}[label=(\arabic*),leftmargin=*]
\item $X_\Sigma$ 在 $\operatorname{Spec} k$ 上紧合当且仅当 $\Sigma$ 是完备扇（见定义 1.7）。
\item $X_\Sigma$ 非异当且仅当每个锥 $\sigma\in\Sigma$ 是标准的，即存在 $M$ 的一组基 $e_1,\ldots,e_n$，使得 $\sigma$ 由 $e_1,\ldots,e_p$ 生成，其中 $p\le n$。
\end{enumerate}

% [原文 p.94]
\begin{example}
设 $\Sigma$ 是例 1.14 中描述的 $\PP^2$ 的扇。该扇有七个锥：$0,\rho_0,\rho_1,\rho_2$ 以及 $\sigma_{0,1},\sigma_{1,2},\sigma_{2,0}$，如第 1 章的图 14 所示。对偶锥为
\begin{gather*}
0^{\vee}=N_\R\\
\rho_0^{\vee}=\{(a,b)\mid a+b\le 0\}\\
\rho_1^{\vee}=\{(a,b)\mid a\ge 0\}\\
\rho_2^{\vee}=\{(a,b)\mid b\ge 0\}\\
\sigma_{0,1}^{\vee}=\{(a,b)\mid a+b\le 0,\ a\ge 0\}\\
\sigma_{1,2}^{\vee}=\{(a,b)\mid a\ge 0,\ b\ge 0\}\\
\sigma_{2,0}^{\vee}=\{(a,b)\mid a+b\le 0,\ b\ge 0\}.
\end{gather*}
注意对每个 $i$（下标模 $3$）都有 $\sigma_{i,i+1}^{\vee}\cap N\cong\N^2$，由此可知 $X_\Sigma$ 被三份 $\mathbb{A}^2_k$ 覆盖。为看出 $X_\Sigma$ 就是 $\PP^2$，我们定义与粘接映射相容的映射 $X_{\sigma_{i,i+1}}\to\PP^2$。特别地，若 $\PP^2$ 有齐次坐标 $x_0,x_1,x_2$，而 $U_i$ 是 $x_i\ne 0$ 的标准仿射子集，则通过如下等同把 $X_{\sigma_{i,i+1}}$ 与 $U_{i+2}$ 等同：
\begin{gather*}
k[\sigma_{0,1}^{\vee}\cap N]=k[x_0/x_2,x_1/x_2],\qquad \text{且 } z^{(0,-1)}=x_0/x_2,\ z^{(1,-1)}=x_1/x_2,\\
k[\sigma_{1,2}^{\vee}\cap N]=k[x_1/x_0,x_2/x_0],\qquad \text{且 } z^{(1,0)}=x_1/x_0,\ z^{(0,1)}=x_2/x_0,\\
k[\sigma_{2,0}^{\vee}\cap N]=k[x_2/x_1,x_0/x_1],\qquad \text{且 } z^{(-1,1)}=x_2/x_1,\ z^{(-1,0)}=x_0/x_1.
\end{gather*}
容易验证这些等同与包含关系 $X_{\rho_i}\subseteq X_{\sigma_{i-1,i}},X_{\sigma_{i,i+1}}$ 相容。例如，我们可以把 $k[\rho_1^{\vee}\cap N]$ 与 $k[x_1/x_0,x_2/x_0,x_0/x_2]$ 等同，其中 $z^{(1,0)}=x_1/x_0$，$z^{(0,1)}=x_2/x_0$，$z^{(0,-1)}=x_0/x_2$。于是显然的包含
\[
  k[x_0/x_2,x_1/x_2]\subseteq k[x_1/x_0,x_2/x_0,x_0/x_2]\supseteq k[x_1/x_0,x_2/x_0]
\]
在上述等同下，恰好与
\[
  k[\sigma_{0,1}^{\vee}\cap N]\subseteq k[\rho_1^{\vee}\cap N]\supseteq k[\sigma_{1,2}^{\vee}\cap N]
\]
重合。这表明 $X_\Sigma$ 同构于 $\PP^2$。\hfill$\square$
\end{example}

由扇构造环面簇在下述意义下是函子性的。设 $M_1$ 与 $M_2$ 是格，并有群同态 $\varphi:M_1\to M_2$。若 $\Sigma_1$ 是 $M_1\otimes\R$ 中的扇、$\Sigma_2$ 是 $M_2\otimes\R$ 中的扇，则称 $\varphi$ 是扇的映射，如果对每个 $\sigma_1\in\Sigma_1$，都存在 $\sigma_2\in\Sigma_2$ 使得 $\varphi(\sigma_1)\subseteq\sigma_2$。这意味着转置映射 ${}^t\varphi:N_2\to N_1$ 满足 ${}^t\varphi(\sigma_2^{\vee})\subseteq\sigma_1^{\vee}$。于是 ${}^t\varphi$ 诱导幺半群环的映射 $k[\sigma_2^{\vee}\cap N_2]\to k[\sigma_1^{\vee}\cap N_1]$，从而诱导映射 $X_{\sigma_1}\to X_{\sigma_2}$。这些映射对 $\sigma_1$ 的各种选择是相容的，因此粘接为态射
\[
  \varphi:X_{\Sigma_1}\to X_{\Sigma_2}.
\]

若诱导映射 $\varphi_\R:M_1\otimes_\Z\R\to M_2\otimes_\Z\R$ 满足 $\varphi_\R^{-1}(|\Sigma_2|)=|\Sigma_1|$，则该态射紧合。

若该态射在对应于锥 $0\in\Sigma_1$ 与 $0\in\Sigma_2$ 的仿射子集上诱导同构，则它是双有理的；这只有在 $\varphi:M_1\to M_2$ 是同构时才会发生。因此 $\varphi$ 是紧合双有理态射，当且仅当 $\varphi$ 是同构且 $\Sigma_1$ 是扇 $\Sigma_2$ 的加细。

% [原文 p.95]
\begin{example}
设 $\sigma\subseteq M_\R\cong\R^n$ 是由一组基 $e_1,\ldots,e_n$ 生成的锥。这给出仿射环面簇 $X_\sigma\cong\mathbb{A}^n_k$。现在通过考虑由
\[
  \Sigma=\{\tau\mid\tau \text{ 是 } \sigma \text{ 的真面}\}\cup\{\tau+\R_{\ge 0}(e_1+\cdots+e_n)\mid\tau \text{ 是 } \sigma \text{ 的真面}\}
\]
给出的扇 $\Sigma$ 来细分锥 $\sigma$。我们在扇中引进了一条新射线，它由 $e_1+\cdots+e_n$ 生成。于是有双有理态射 $X_\Sigma\to X_\sigma$。作为练习，证明该态射是 $X_\sigma$ 在原点处的爆破。
\end{example}

格 $M$ 与 $N$ 在环面簇的语境中有如下解释。首先，由于扇总含有锥 $0$，而 $X_0=\operatorname{Spec} k[N]$ 是代数环面，故 $X_\Sigma$ 总含有同构于代数环面的开子集。该子集常称为 $X_\Sigma$ 中的大环面，或大环面轨道。

一般地，若 $L$ 是格，我们用记号
\[
  \mathbb{G}(L):=\operatorname{Spec} k[\operatorname{Hom}_\Z(L,\Z)]
\]
表示由 $L$ 确定的代数环面。$\mathbb{G}(L)$ 的 $k$-值点集是 $L\otimes_\Z k^{\times}$。

现在环面 $\mathbb{G}(M)$ 作用在 $X_\Sigma$ 上。该作用 $\mathbb{G}(M)\times_k X_\Sigma\to X_\Sigma$ 在仿射集 $X_\sigma$ 上由环映射
\[
  k[\sigma^{\vee}\cap N]\to k[N]\otimes_k k[\sigma^{\vee}\cap N]
\]
给出，该映射为
\[
  z^n\mapsto z^n\otimes z^n.
\]
这个作用显然在不同开集 $X_\sigma$ 上相容，于是得到 $X_\Sigma$ 上的环面作用。

扇 $\Sigma$ 的锥实际上与 $\mathbb{G}(M)$ 在 $X_\Sigma$ 上的轨道一一对应。事实上，对应于 $\tau\in\Sigma$ 的环面轨道是
\[
  X_\tau\setminus\bigcup_{\omega\subsetneq\tau}X_\omega,
\]
其中并取遍 $\tau$ 的所有真面。对该轨道中的点 $x$，子环面 $\mathbb{G}(\R\tau\cap M)\subseteq\mathbb{G}(M)$ 是 $x$ 的稳定化子。我们把这个轨道的闭包记作 $D_\tau$。对应 $\tau\mapsto D_\tau$ 是反包含的。我们把子簇 $D_\tau$（$\tau\in\Sigma$）称为 $X_\Sigma$ 的环面层。

$D_\tau$ 本身是环面簇，其扇由商扇 $\Sigma(\tau)$ 给出（见定义 1.10）。注意若 $\rho$ 是射线，则 $D_\rho$ 是 $X_\Sigma$ 上的（Weil）除子。

再来看 $N$。它可以解释为代数环面 $\mathbb{G}(M)$ 的特征格，即 $N=\operatorname{Hom}(\mathbb{G}(M),\mathbb{G}_m)$。这里 $\mathbb{G}_m=\mathbb{G}(\Z)$ 是乘法群。元素 $n\in N$ 给出映射 $\Z\to N$，从而给出映射 $k[\Z]\to k[N]$，得到特征 $\mathbb{G}(M)\to\mathbb{G}_m$。这个特征通常写作 $z^n$，可以看作正则函数，它在 $\mathbb{G}_m(M)$ 的 $k$-值点集 $M\otimes_\Z k^{\times}$ 上由 $m\otimes z\mapsto z^{\langle n,m\rangle}$ 给出。

% [Tutor Note: p.95 上一行的 $\mathbb{G}_m(M)$ 应为 $\mathbb{G}(M)$：同段开头已写 $N=\operatorname{Hom}(\mathbb{G}(M),\mathbb{G}_m)$，而 $\mathbb{G}_m=\mathbb{G}(\Z)$ 只是一维乘法群，其 $k$-值点集是 $k^\times$ 而非 $M\otimes_\Z k^\times$；且下文一律用 $\mathbb{G}(M)$。此处为原文笔误，照译以存真。]

给定扇 $\Sigma$，$z^n$ 是 $\mathbb{G}(M)\subseteq X_\Sigma$ 上的正则函数，因而定义了 $X_\Sigma$ 上的有理函数。由于环面簇总是正规的，谈论 $z^n$ 在 Weil 除子 $D_\rho$ 上的消没阶是有意义的，其中 $\rho\in\Sigma[1]$，而 $\Sigma[1]$ 表示 $\Sigma$ 中一维锥的集合。这个零点阶实际上就是 $\langle n,m_\rho\rangle$，其中 $m_\rho$ 是射线 $\rho$ 的本原整生成元。于是
\[
  \sum_{\rho\in\Sigma[1]}\langle n,m_\rho\rangle D_\rho
\]
是主除子。

% [原文 p.96]
可以证明 $X_\Sigma$ 的 Weil 除子类群如下给出。设 $T_\Sigma$ 表示如 §1.3 中由 $\Sigma[1]$ 生成的自由 Abel 群，对每个 $\rho\in\Sigma[1]$ 有生成元 $t_\rho$。回忆映射 $r:T_\Sigma\to M$，它由 $r(t_\rho)=m_\rho$ 定义，其中 $m_\rho$ 是 $\rho$ 的本原生成元。转置映射 ${}^tr:N\to T_\Sigma^{\vee}$ 由
\[
  n\mapsto\sum_{\rho\in\Sigma[1]}\langle n,m_\rho\rangle t_\rho^{*}
\]
给出。则 $\operatorname{Cl}(X_\Sigma)\cong\operatorname{coker}{}^tr$。这里，$T_\Sigma^{\vee}$ 中的元素 $\psi:T_\Sigma\to\Z$ 对应于 Weil 除子 $\sum_{\rho\in\Sigma[1]}\psi(t_\rho)D_\rho$，而 ${}^tr(n)$ 是 $z^n$ 的零点与极点除子。

若 $r$ 实际上是满射——例如当 $X_\Sigma$ 非异且在 $\operatorname{Spec} k$ 上紧合时就是如此——我们得到正合列
\begin{equation}
  0\to K_\Sigma\to T_\Sigma\to M\to 0. \tag{3.1}
\end{equation}
于是对偶正合列为
\begin{equation}
  0\to N\to T_\Sigma^{\vee}\to\operatorname{Cl}(X_\Sigma)\to 0, \tag{3.2}
\end{equation}
由 $\psi\in T_\Sigma^{\vee}$ 诱导的除子是 Cartier 除子，如果 $\psi$ 由分片线性函数 $\varphi:|\Sigma|\to\R$ 诱导，且 $\varphi(m_\rho)=\psi(t_\rho)$。

给定由分片线性函数 $\varphi$ 定义的 Cartier 除子 $D$，$D$ 是非常丰沛的，当且仅当 $\varphi$ 严格凸。

当然，Cartier 除子定义了线丛。给定由分片线性函数 $\varphi$ 指定的 Cartier 除子，用平凡化很容易描述这个线丛 $\cL_\varphi$。对每个极大锥 $\sigma\in\Sigma$，$\varphi|_\sigma$ 由某个 $n_\sigma\in N$ 给出。于是 $\cL_\varphi$ 自然地与平凡线丛 $\cO_{X_\sigma}\cdot z^{-n_\sigma}$ 等同。使用这个平凡化，从 $X_\sigma$ 到 $X_{\sigma'}$ 的转移映射由乘以 $z^{n_{\sigma'}-n_\sigma}$ 给出。

若 $\tau\in\Sigma$，则线丛 $\cL_\varphi$ 到环面层 $D_\tau$ 的限制容易描述：它由扇 $\Sigma(\tau)$ 上的分片线性函数 $\varphi(\tau)$ 给出（见定义 1.10）。这在相差一个线性函数的意义下是良定义的，而这已足以指定线丛 $\cL_\varphi|_{D_\tau}$。

% [Tutor Note: p.96 (3.1)(3.2) 的核对锚点：取基 $t_\rho$ 与 $n_i$，${}^tr$ 的矩阵元为 $\langle n_i,m_\rho\rangle$，故 (3.2) 正是 (3.1) 取 $\operatorname{Hom}(-,\Z)$ 的对偶；$\operatorname{Cl}(X_\Sigma)\cong\operatorname{coker}{}^tr$ 与“除子类群 $=T_\Sigma^\vee/\operatorname{im}{}^tr$”一致，可用来检查式 (3.2) 中两端的次序。]

\subsubsection{来自多面体的环面簇}

在环面簇射影于仿射环面簇的情形，实际上有一个非常重要的对偶观点。

设 $\Delta\subseteq N_\R$ 是至少有一个顶点的格多面体。我们可以如下定义一个簇。首先，令 $C(\Delta)\subseteq N_\R\oplus\R$ 由
\[
  C(\Delta)=\overline{\{(rn,r)\mid n\in\Delta,\ r\ge 0\}}
\]
定义。这是 $\Delta$ 上的锥，上划线表示闭包。该锥是有理多面体的，且
\[
  C(\Delta)\cap(N_\R\oplus\{0\})=\operatorname{Asym}(\Delta),
\]
其中 $\operatorname{Asym}(\Delta)$ 是 $\Delta$ 的渐近锥，定义为 $r\Delta$ 当 $r\to 0$ 时的 Hausdorff 极限。若 $\Delta$ 紧，则当然 $\operatorname{Asym}(\Delta)=\{0\}$。

% [Tutor Note: p.96 渐近锥“$r\Delta$ 当 $r\to 0$ 时的 Hausdorff 极限”应为 $r\to\infty$（或等价地 $(1/r)\Delta$ 当 $r\to 0$）：若 $r\to 0$，集合 $r\Delta$ 收缩到原点，极限只能是 $\{0\}$，与紧情形 $\operatorname{Asym}(\Delta)=\{0\}$ 以外的用法矛盾。此为原文笔误，照译以存真。]

现在 $k[C(\Delta)\cap(N\oplus\Z)]$ 是有限生成的 $k$-代数，它带有由投影 $N\oplus\Z\to\Z$ 给出的自然分次，即 $\deg z^{(n,d)}=d$。该环的零次部分是 $k[\operatorname{Asym}(\Delta)\cap N]$，于是得到簇
\[
  \PP_\Delta:=\operatorname{Proj}k[C(\Delta)\cap(N\oplus\Z)],
\]
它在 $\operatorname{Spec} k[\operatorname{Asym}(\Delta)\cap N]$ 上射影。

% [原文 p.97]
$\PP_\Delta$ 实际上是环面簇：它同构于 $X_{\check\Sigma_\Delta}$，其中 $\check\Sigma_\Delta$ 是如定义 1.9 所定义的 $\Delta$ 的法扇。

\begin{example}
设 $\Delta=\operatorname{Conv}\{(0,0),(1,0),(0,1)\}$。则 $k[C(\Delta)\cap(N\oplus\Z)]\cong k[x_0,x_1,x_2]$，其中 $z^{(0,0,1)}\mapsto x_0$，$z^{(1,0,1)}\mapsto x_1$，$z^{(0,1,1)}\mapsto x_2$。分次是标准的。因此 $\PP_\Delta\cong\PP^2$。注意 $\Delta$ 的法扇正是例 3.2 中给出的 $\PP^2$ 的扇。更一般地，标准 $n$-单形给出 $\PP^n$。
\end{example}

用这种方式以多面体定义环面簇，能得到更几何的图景。$\Delta$ 的面与 $\Delta$ 的法扇中的锥之间存在一一的反包含对应，因此 $\Delta$ 的面与 $\PP_\Delta$ 的环面层之间存在一一对应，而后者现在是保包含的。于是可以把 $\Delta$ 的面看成以几何方式指示各层如何相交的组合数据。此外，若 $\sigma$ 是 $\Delta$ 的面，则 $D_{\check\sigma}\cong\PP_\sigma$，其中 $\check\sigma=N_\Delta(\sigma)$ 是 $\Delta$ 沿 $\sigma$ 的法锥。因此从这个图景很容易读出环面簇及其环面层的几何。

\begin{remark}
我们在此记录几个基本事实。首先，由于 $\PP_\Delta$ 定义为 Proj，它自然带有线丛 $\cO_{\PP_\Delta}(1)$。该线丛由如定义 1.9 中定义的分片线性函数 $\varphi_\Delta:|\check\Sigma_\Delta|\to\R$ 确定。$\Gamma(\PP_\Delta,\cO_{\PP_\Delta}(1))$ 有一组由 $\Delta$ 的整点集（即 $\Delta\cap N$）索引的基。这一点是清楚的，因为这些整点索引了 $k[C(\Delta)\cap(N\oplus\Z)]$ 的一次齐次部分的一组基。其次，面 $\sigma\subseteq\Delta$ 确定了 $\Delta$ 的法扇中的锥 $N_\Delta(\sigma)$，而由 (1.3)，该锥的对偶恰好是切楔 $T_\sigma\Delta$。因此面 $\sigma\subseteq\Delta$ 在确定 $\PP_\Delta$ 的一个闭的环面层的同时，也确定了 $\PP_\Delta$ 的开仿射子集 $\operatorname{Spec} k[(T_\sigma\Delta)\cap N]$。

这为平凡化线丛 $\cO_{\PP_\Delta}(1)$ 提供了方便的描述。对 $\Delta$ 的每个顶点 $v$，我们可以在 $\operatorname{Spec} k[(T_v\Delta)\cap N]$ 上平凡化 $\cO_{\PP_\Delta}(1)$，使得对应于某个 $n\in\Delta\cap N$ 的截面变成单项式 $z^{n-v}$。顶点 $v$ 与 $v'$ 所对应开集上的平凡化之间的转移映射于是由乘以 $z^{v-v'}$ 给出。
\end{remark}

下面这个例子是本书的关键。

\begin{example}
\textit{Mumford 退化。}考虑下述情形，它已在 §1.1 中出现过。设 $\Delta\subseteq N_\R$ 是紧格多面体，$\cP$ 是把 $\Delta$ 分解为格多面体的一个多面体分解，$\varphi:\Delta\to\R$ 是具有整斜率的凸分片线性函数。于是与 §1.1 类似，有多面体
\[
  \tilde\Delta=\{(n,r)\in N_\R\oplus\R\mid n\in\Delta,\ r\ge\varphi(n)\}.
\]
$\tilde\Delta$ 的渐近锥是 $0\times\R_{\ge 0}\subseteq N_\R\times\R$，故 $k[C(\tilde\Delta)\cap(N\oplus\Z\oplus\Z)]$ 是 $k[\N]$-代数。环面簇 $\PP_{\tilde\Delta}$ 于是带有一个射影态射 $\pi:\PP_{\tilde\Delta}\to\mathbb{A}^1_k$。

为更清楚地了解 $\pi$ 为我们做了什么，先描述 $\tilde\Delta$ 的法扇 $\check\Sigma_{\tilde\Delta}$。注意 $\tilde\Delta$ 有两类面。若 $p:N_\R\oplus\R\to N_\R$ 表示投影，则 $\tilde\Delta$ 有竖直面，它们经 $p$ 非同胚地投影到 $\Delta$ 的面；还有 $\tilde\Delta$ 的水平面，它们同胚地投影到 $\cP$ 的元素。若 $\sigma$ 是竖直面，则法锥 $\check\sigma:=N_{\tilde\Delta}(\sigma)$ 位于 $M_\R\times\{0\}$ 中，并且事实上是 $\Delta$ 的法扇中的锥。另一方面，对极大水平面 $\sigma$（此时 $\varphi|_{p(\sigma)}$ 的斜率为 $m_\sigma\in M$），其法锥是由 $(-m_\sigma,1)$ 生成的射线。
\end{example}

% [原文 p.98]
\begin{figure}[htbp]
\centering
\includegraphics[width=0.32\textwidth]{f3_1.png}
\figtitle{1}{}
\end{figure}

现在态射 $\pi:\PP_{\tilde\Delta}\to\mathbb{A}^1_k$ 当然是正则函数，它由单项式 $z^\rho$ 定义，其中 $\rho=(0,1)\in\operatorname{Asym}(\tilde\Delta)\subseteq N_\R\oplus\R$。由于 $\check\Sigma_{\tilde\Delta}$ 的射线的本原生成元要么形如 $(m,0)$，要么对某个 $m\in M$ 形如 $(m,1)$，故 $z^\rho$ 在对应于第一类射线的除子上不消没，而在对应于第二类射线的除子上恰消没到一阶。因此 $\pi^{-1}(0)$ 同构于 $\PP_{\tilde\Delta}$ 中对应于余维一水平面的环面除子之并。这些除子与 $\cP_{\max}$ 的元素一一对应，我们可以写
\[
  \pi^{-1}(0)=\bigcup_{\sigma\in\cP_{\max}}\PP_\sigma.
\]
当然，这些环面簇恰按 $\cP$ 所指示的方式相交。例如，设 $\Delta$ 与 $\cP$ 如图 1 所示。则 $\pi^{-1}(0)$ 是三个 $\PP^2$ 和一个 $\PP^1\times\PP^1$ 的并。

其次，注意 $\PP_{\tilde\Delta}\setminus\pi^{-1}(0)$ 同构于 $\PP_\Delta\times\mathbb{G}_m$。事实上，只需把环 $k[C(\tilde\Delta)\cap(N\oplus\Z\oplus\Z)]$ 在元素 $z^{(0,1,0)}$ 处局部化。这与把 $\tilde\Delta$ 换成 $\Delta\times\R$ 是一回事，而 $\PP_{\Delta\times\R}=\PP_\Delta\times_k\operatorname{Spec} k[\Z]=\PP_\Delta\times_k\mathbb{G}_m$。这说明 $\PP_{\tilde\Delta}\setminus\pi^{-1}(0)\cong\PP_\Delta\times_k\mathbb{G}_m$。

综合以上，我们看到 $\pi$ 是环面簇的退化，其中对 $t\ne 0$ 有 $\pi^{-1}(t)\cong\PP_\Delta$，而 $\pi^{-1}(0)$ 是环面簇之并，其相交由 $\cP$ 描述。

\subsection{对数概形}

我们直接从定义开始。

\begin{definition}
概形 $X$ 上的前对数结构是 $X$ 上的幺半群层 $\cM_X$，连同幺半群层的同态
\[
  \alpha_X:\cM_X\to\cO_X,
\]
其中 $\cO_X$ 上的幺半群结构由乘法给出。

称前对数结构为对数结构，如果
\[
  \alpha_X:\alpha_X^{-1}(\cO_X^{\times})\to\cO_X^{\times}
\]
是同构。这里 $\cO_X^{\times}$ 表示 $\cO_X$ 的可逆元层。

对数概形是带有对数结构 $\alpha_X:\cM_X\to\cO_X$ 的概形 $X$。我们通常把对数概形写作 $X^\dagger$，而把 $\cM_X,\alpha_X$ 隐含在记号中。
\end{definition}

% [原文 p.99]
对数概形的态射 $f:X^\dagger\to Y^\dagger$ 是概形的态射 $f:X\to Y$，连同幺半群层的同态 $f^{\#}:f^{-1}\cM_Y\to\cM_X$，使得图表
\[
  \begin{array}{ccc}
    f^{-1}\cM_Y & \xrightarrow{f^{\#}} & \cM_X\\[10pt]
    {\scriptstyle\alpha_Y}\big\downarrow & & \big\downarrow{\scriptstyle\alpha_X}\\[10pt]
    f^{-1}\cO_Y & \xrightarrow[f^{*}]{} & \cO_X
  \end{array}
\]
交换。这里 $f^*$ 是由态射 $f$ 定义的通常的正则函数拉回。

这些定义虽然直截了当，却是我一直觉得难以真正吸收和消化的，学会以对数结构的方式思考需要一些时间。我将通过给出大量例子来帮助这一过程。

不过在进入这些例子之前，有几点需要说明。首先，我们应简要提及这个定义与前面关于对数微分的讨论之间的联系。本质上，想法是幺半群层 $\cM_X$ 指定了“我们可以取 $d\log$ 的对象”。我们将在 §3.3 讨论对数微分时使之精确化。

其次，我们尚未指定幺半群层 $\cM_X$ 生活在哪种拓扑中。一般地，处理概形时应当在 étale 拓扑而非 Zariski 拓扑中工作，因为确实存在重要的对数结构例子，若只考虑 Zariski 开子集就会丢失重要信息；见例 3.14。若改为在复解析空间上工作，则总可以舒适地在解析拓扑中工作，对 étale 拓扑这一概念感到不适的读者最好在解析拓扑中思考问题。尽管如此，对本书中将出现的大多数对数结构而言，Zariski 拓扑已经足够。

然而，由于我们正式将在 étale 拓扑中工作，值得为不熟悉它的读者作几点说明。更多细节可在 Milne 的书 [83] 中找到。在 étale 拓扑中，人们用 étale 态射 $U\to X$ 代替开集 $U\subseteq X$ 的概念。点 $x\in X$ 的 étale 邻域是像包含 $x$ 的 étale 映射 $U\to X$。étale 拓扑中的层 $\cF$ 于是对每个 étale 态射 $U\to X$ 指定一个 Abel 群（或幺半群等）$\cF(U)$。它们照常带有限制映射：给定 $U\to V\to X$，有映射 $\cF(V)\to\cF(U)$。然后要求 $\cF$ 满足通常的层公理（适当陈述，例如见 [83]，第 II 章 §1）。

当我们谈论层的茎时，需要选取概形 $X$ 的一个几何点。设 $x\in X$ 有剩余域 $k(x)$，并固定 $k(x)$ 的一个可分闭包 $\overline{k(x)}$。令 $\bar x=\operatorname{Spec}\overline{k(x)}$，并令 $\bar x\to X$ 为诱导映射。这称为 $X$ 的几何点。然后我们定义 $\cF$ 在 $\bar x$ 处的茎 $\cF_{\bar x}$ 为群（幺半群等）$\cF(U)$ 沿图表
\[
  \begin{array}{ccc}
    \bar x & \longrightarrow & U\\[6pt]
    \searrow & & \big\downarrow\\[6pt]
    & & X
  \end{array}
\]
的直接极限。

% [原文 p.100]
例如，étale 拓扑中的层 $\cO_X$ 由 $\cO_X(U)=\Gamma(U,\cO_U)$ 定义，其中 $\cO_U$ 是 $U$ 上通常的结构层。茎 $\cO_{X,\bar x}$ 比通常的局部环 $\cO_{X,x}$ 大得多。

这些细节在本书中对我们不会起特别重要的作用，但似乎有必要引入这些记号，以使此处陈述的命题在技术上是正确的。

\begin{examples}
\begin{enumerate}[label=(\arabic*),start=0,leftmargin=*]
\item 若 $X$ 是概形，$\cM_X=\cO_X^{\times}$，$\alpha_X$ 为包含映射，则得到 $X$ 上的一个对数结构，称为平凡对数结构。带有平凡对数结构的概形之间的对数态射与概形态射是一回事。
\item \textit{基本例子。}设 $X$ 是概形，$D\subseteq X$ 是纯余维一的闭子集。令 $j:X\setminus D\hookrightarrow X$ 为包含映射，并考虑
\[
  \cM_{(X,D)}:=(j_*\cO_{X\setminus D}^{\times})\cap\cO_X.
\]
这是 $X$ 上在 $X\setminus D$ 上可逆的正则函数层。我们取 $\alpha_X:\cM_{(X,D)}\hookrightarrow\cO_X$ 为包含映射。这显然是对数结构，称为由 $D$ 诱导的除子对数结构。
\item 设 $k$ 是域，$X=\operatorname{Spec} k$。令 $\cM_X=k^{\times}\oplus Q$，其中 $Q$ 是唯一可逆元为 $0\in Q$ 的幺半群。令 $\alpha_X:k^{\times}\oplus Q\to k$ 由
\[
  \alpha_X(x,q):=
  \begin{cases}
    x & q=0,\\
    0 & q\ne 0
  \end{cases}
\]
给出。这是对数结构。有两个特殊情形：若 $Q=\{0\}$，则得到 $X$ 上的平凡对数结构。若 $Q=\N$，则得到所谓的标准对数点。我们通常把标准对数点写作 $\operatorname{Spec} k^\dagger$。
\end{enumerate}
\end{examples}

除子对数结构与标准对数点之间存在一种联系，它来自下述一般构造。

\begin{definition}
设 $\alpha:P_X\to\cO_X$ 是 $X$ 上的前对数结构。则该前对数结构相伴的对数结构由幺半群层
\[
  \cM_X:=\frac{P_X\oplus\cO_X^{\times}}{\{(p,\alpha(p)^{-1})\mid p\in\alpha^{-1}(\cO_X^{\times})\}}
\]
给出，而 $\alpha_X:\cM_X\to\cO_X$ 由
\[
  \alpha_X(p,f):=\alpha(p)\cdot f
\]
给出。这显然是对数结构。
\end{definition}

\begin{definition}
若 $f:X\to Y$ 是概形的态射，且 $Y$ 带有对数结构 $\alpha_Y:\cM_Y\to\cO_Y$，则 $X$ 上的拉回对数结构是由复合
\[
  \alpha:f^{-1}(\cM_Y)\xrightarrow{\alpha_Y}f^{-1}(\cO_Y)\xrightarrow{f^*}\cO_X
\]
给出的前对数结构所相伴的对数结构。我们把该拉回对数结构的幺半群层写作 $f^*\cM_Y$。
\end{definition}

\begin{example}
设 $X=\operatorname{Spec} k$，$Y=\operatorname{Spec} k[x]=\mathbb{A}^1_k$，并令 $f:X\to Y$ 把 $X$ 映到 $Y$ 的原点。考虑由 $\{0\}\subseteq Y$ 给出的 $\mathbb{A}^1_k$ 上的除子对数结构，并经 $f$ 把它拉回到 $X$。此时 $f^{-1}\cM_Y$ 就是层 $\cM_Y$ 在 $0$ 处的茎；

% [原文 p.101]
它由形如 $\varphi\cdot x^n$（$n\ge 0$）的函数芽组成，其中 $\varphi$ 是在 $0$ 的（étale）邻域中定义的可逆函数。映射
\[
  \alpha:f^{-1}\cM_Y\to f^{-1}\cO_Y\to\cO_X=k
\]
是 $\varphi\cdot x^n$ 在 $x=0$ 处的取值，因此
\[
  \alpha(\varphi\cdot x^n)=
  \begin{cases}
    \varphi(0) & n=0,\\
    0 & n>0.
  \end{cases}
\]
这定义了 $X$ 上的前对数结构。取相伴的对数结构，我们考虑
\[
  \frac{f^{-1}\cM_Y\oplus k^{\times}}{\{(\varphi,\varphi(0)^{-1})\mid \varphi\in\cO_{Y,\bar 0}^{\times}\}}.
\]
它实际上通过映射
\[
  f^{-1}\cM_Y\oplus k^{\times}\ni(\varphi\cdot x^n,s)\mapsto(\varphi(0)\cdot s,n)
\]
同构于 $k^{\times}\oplus\N$。此外
\[
  \alpha_X(\varphi\cdot x^n,s)=
  \begin{cases}
    \varphi(0)\cdot s & n=0,\\
    0 & n>0,
  \end{cases}
\]
因此 $X$ 上的拉回对数结构给出标准对数点。\hfill$\square$
\end{example}

通过考虑正合列
\[
  1\to\cO_X^{\times}\xrightarrow{\alpha_X^{-1}}\cM_X\to\overline{\cM}_X\to 0
\]
来理解对数结构，常常是有用的。加法地书写的幺半群层 $\overline{\cM}_X$ 称为 $X^\dagger$ 的幽灵层，应当把它看作包含了对数结构的组合信息。

% [Tutor Note: p.101 幽灵层：$\overline{\cM}_X$ 是乘法记法下的商层 $\cM_X/\cO_X^\times$，原文说的“按加法书写”指把它作为加法幺半群层；上面的正合列在群层意义下 $1\to\cO_X^\times\to\cM_X\to\overline{\cM}_X\to 0$ 成立。核对时可与例 3.12 中的映射 $q$ 对照：$q$ 正是穿过 $\overline{\cM}_{(X,D),\bar x}$ 的消没阶映射。]

\begin{example}
设 $X$ 局部 Noether 且正规，$D\subseteq X$ 是纯余维一的闭子集。则对 $x\in X$，令 $r$ 为 $D$ 的含 $x$ 的不可约分支个数。有映射
\[
  q:\cM_{(X,D),\bar x}\to\N^r,
\]
它把正则函数 $f\in\cM_{(X,D),\bar x}$ 对应到 $f$ 沿 $D$ 的含 $x$ 的 $r$ 个分支的消没阶。该映射显然经由 $\overline{\cM}_{(X,D),\bar x}$ 分解。另一方面，若 $f_1,f_2\in\cM_{(X,D),\bar x}\subseteq\cO_{X,\bar x}$ 沿 $D$ 有相同的消没阶，则 $f_1f_2^{-1}$ 在 $\bar x$ 的邻域内每个素除子上的消没阶为零。由于 $X$ 正规，由此推出（例如见 [57] 第 132 页给出的论证）$f_1=f_2\cdot h$，其中 $h\in\cO_{X,\bar x}^{\times}$。因此 $\overline{\cM}_{(X,D),\bar x}$ 是 $\N^r$ 的子幺半群。\hfill$\square$
\end{example}

注意若 $f:X\to Y$ 是概形的态射，且 $Y$ 带有在 $X$ 上给出拉回对数结构的对数结构，则 $\overline{\cM}_X=f^{-1}\overline{\cM}_Y$。事实上，我们有

% [原文 p.102]
图表
\[
  \begin{array}{c@{\quad}c@{\quad}c@{\quad}c@{\quad}c@{\quad}c@{\quad}c@{\quad}c}
    & & 1 & & 1 & & & \\[2pt]
    & & \big\downarrow & & \big\downarrow & & & \\[2pt]
    & & K & \xrightarrow{\hspace{1.2em}=\hspace{1.2em}} & K & & & \\[2pt]
    & & \big\downarrow & & \big\downarrow & & & \\[2pt]
    1 & \longrightarrow & f^{-1}\cO_Y^{\times}\oplus\cO_X^{\times} & \longrightarrow & f^{-1}\cM_Y\oplus\cO_X^{\times} & \longrightarrow & f^{-1}\overline{\cM}_Y & \longrightarrow\ 0\\[2pt]
    & & \big\downarrow & & \big\downarrow & & \overset{=}{\big\downarrow} & \\[2pt]
    1 & \longrightarrow & \cO_X^{\times} & \longrightarrow & \cM_X & \longrightarrow & \overline{\cM}_X & \longrightarrow\ 0\\[2pt]
    & & \big\downarrow & & \big\downarrow & & & \\[2pt]
    & & 1 & & 1 & & &
  \end{array}
\]
其中 $K$ 既是映射 $f^{-1}\cO_Y^{\times}\oplus\cO_X^{\times}\to\cO_X^{\times}$（由 $(\varphi_1,\varphi_2)\mapsto f^*(\varphi_1)\cdot\varphi_2$ 给出）的核，也是映射 $f^{-1}\cM_Y\oplus\cO_X^{\times}\to\cM_X$ 的核。（幺半群态射的核一般不存在，但若态射是“按子群取商”给出的则存在。）后一个核按定义由
\[
  \{(\varphi_1,\varphi_2)\in f^{-1}\cM_Y\oplus\cO_X^{\times}\mid f^*\alpha_Y(\varphi_1)=\varphi_2^{-1}\}
\]
给出。容易验证蛇引理在此语境下成立。

% [Tutor Note: p.102 蛇引理图：幺半群范畴没有 Abel 结构，通常意义的核/余核不存在；原文括注“若态射按子群取商给出”是关键——本图中两个水平正合列与竖列均为群层（或可差核）意义下正合，故蛇引理可用。逐列核对时注意右列竖直方向是 $=$ 而非箭头，表示 $\overline{\cM}_Y\to\overline{\cM}_X$ 为同构（由上一段 $\overline{\cM}_X=f^{-1}\overline{\cM}_Y$ 得出）。]

\begin{example}
设 $X=\mathbb{A}^2_k=\operatorname{Spec} k[x,y]$，并令 $D\subseteq X$ 由 $xy=0$ 给出，即 $D=V(xy)$。于是得到 $X$ 上的除子对数结构 $\cM_{(X,D)}$，而 $\overline{\cM}_{(X,D)}$ 是层 $(i_{1*}\N)\oplus(i_{2*}\N)$，其中 $i_1:V(x)\hookrightarrow X$ 与 $i_2:V(y)\hookrightarrow X$ 是两条坐标轴的包含映射。映射 $\cM_{(X,D)}\to(i_{1*}\N)\oplus(i_{2*}\N)$ 把开集 $U$ 上在 $U\setminus D$ 上可逆的正则函数 $f$ 映到 $f$ 在 $U\cap V(x)$ 与 $U\cap V(y)$ 上的消没阶。如例 3.12 那样，该映射是满射，其核为 $\cO_X^{\times}$。

若经包含映射 $D\hookrightarrow X$ 拉回 $\cM_{(X,D)}$，则得到 $D$ 上的对数结构，且 $\overline{\cM}_D=i_{1*}\N\oplus i_{2*}\N$（此时把 $i_1,i_2$ 看作坐标轴到 $D$ 的包含映射）。特别地，$\cM_D$ 的某些截面不能看作 $D$ 上的函数。应当把 $D$ 上的对数结构看作在记忆 $D$ 如何位于 $X$ 之中的某些信息。
\end{example}

\begin{example}
设我们有如下形状的除子 $D\subseteq X$。$D$ 是不可约的非正规曲面，其正规化 $\tilde D$ 同构于 $E\times\PP^1$，其中 $E$ 是椭圆曲线。设 $\tau\in E$ 是二阶挠点，$P\in\PP^1$ 是一点。我们假定映射 $\tilde D\to D$ 把点 $e\times\{P\}$ 与点 $(e+\tau)\times\{P\}$ 等同。令 $E'$ 是 $E$ 按平移 $\tau$ 取的商。则存在映射 $D\to E'$，其纤维是两个交于一点的 $\PP^1$ 之并。$D$ 的奇点轨迹同构于 $E'$。现在我们有 $X$ 上的除子对数结构 $\cM_{(X,D)}$，利用 $E'$ 与 $D$ 的奇点轨迹的等同，可以把它限制为 $E'$ 上的对数结构 $\cM_{E'}$。注意此时 $\overline{\cM}_{E'}$ 是以 $\N^2$ 为纤维的局部常值层，但只在 étale 拓扑或解析拓扑中如此，在 Zariski 拓扑中不然。在 $E'$ 的任一点处、在以上两种拓扑中任一拓扑下局部地看，$D$ 都有两个不可约分支，但整体上这些分支会被互换。
\end{example}

% [原文 p.103]
所以局部常值层 $\overline{\cM}_{E'}$ 具有非平凡的单值性（monodromy），它互换 $\N^2$ 的生成元。在 Zariski 拓扑中我们看不到这种单值性。

这解释了一般情况下 Zariski 拓扑为何不够用。然而，在本书所需的情形中永远不会出现这种类型的例子。

% [Tutor Note: p.103 “非平凡单值性”的核对锚点：$\overline{\cM}_{E'}$ 在 étale/解析拓扑下局部常值、纤维 $\N^2$，整体上两生成元被 $\tau$-平移互换，故其单值表示 $\pi_1(E')\to\operatorname{GL}_2(\Z)$ 非平凡；Zariski 局部化看不到它，正因为 $D$ 的两个分支无法在 Zariski 开集上分开。]

\begin{example}
我们下一个（篇幅较长的）例子探讨对数态射所携带的数据。我们将看到，一个对数态射可以携带大量额外数据。令 $X^\dagger$ 为例 3.13 中的仿射平面。我们可以如下理解这个例子。照常固定 $M=\Z^2$、$M_\R$、$N$、$N_\R$，并令 $\sigma\subseteq M_\R$ 为第一象限。则 $X=X_\sigma$，因为 $\sigma^{\vee}\cap N=\N^2$。

现在考虑从标准对数点到 $X^\dagger$ 的对数态射 $f:\operatorname{Spec} k^\dagger\to X^\dagger$，其像为原点。与该对数态射相伴的额外信息是一个交换图表
\[
  \begin{array}{ccc}
    \cM_{(X,D),\bar 0} & \xrightarrow{f^{\#}} & k^{\times}\oplus\N\\[10pt]
    {\scriptstyle\alpha_X}\big\downarrow & & \big\downarrow{\scriptstyle\alpha_{\operatorname{Spec}k^{\dagger}}}\\[10pt]
    \cO_{X,\bar 0} & \xrightarrow[f^{*}]{} & k
  \end{array}
\]
这里 $f^*$ 只是把函数芽在 $0$ 处取值，而 $\cM_{(X,D),\bar 0}$ 由在 $X\setminus D$ 上可逆的函数芽组成，即局部地形如 $\varphi\cdot z^n$，其中 $n\in\sigma^{\vee}\cap N$，$\varphi\in\cO_{X,\bar 0}^{\times}$。注意 $f^{\#}$ 诱导映射 $\N^2=\overline{\cM}_{(X,D),\bar 0}\to\cM_{\operatorname{Spec}k^\dagger}=\N$，这是幺半群同态，因此这是由该态射决定的离散数据。该映射由 $n\mapsto\langle n,m\rangle$ 给出，其中 $m\in\sigma\cap M$。此外 $m\in\operatorname{Int}(\sigma)$。事实上，假设 $m=(a,0)$，其中 $a\in\N$。则 $f^{\#}(z^{(0,1)})=(x,0)\in k^{\times}\oplus\N$，其中 $x\in k^{\times}$，而 $x=\alpha_{\operatorname{Spec}k^\dagger}(f^{\#}(z^{(0,1)}))=f^*\alpha_X(z^{(0,1)})=0$，矛盾。类似地，$m=(0,a)$ 也导出矛盾。

还要注意，若 $\varphi$ 可逆，则 $f^*\alpha_X(\varphi)=\varphi(0)$，故 $f^{\#}(\varphi)=(\varphi(0),0)$。由此断定 $f^{\#}$ 完全由 $m$ 以及幺半群映射 $g:\sigma^{\vee}\cap N\to k^{\times}$ 决定，使得
\begin{equation}
  f^{\#}(\varphi z^n)=(\varphi(0)g(n),\langle n,m\rangle). \tag{3.3}
\end{equation}
所以我们看到，该映射完全由 $m\in\operatorname{Int}(\sigma)\cap M$ 与 $g\in M\otimes_\Z k^{\times}$ 确定。

让我们对此稍作展开，以尝试在一定程度上解释这些额外数据的含义。考虑通过引进另一条射线 $\R_{\ge 0}m'$（其中 $m'\in\operatorname{Int}(\sigma)\cap M$ 是本原的）得到的锥 $\sigma$ 的细分。这给出有三条射线 $\rho_1=\R_{\ge 0}(1,0)$、$\rho_2=\R_{\ge 0}(0,1)$、$\rho=\R_{\ge 0}m'$ 以及两个二维锥的扇 $\Sigma$。$M$ 上的恒等映射诱导从 $\Sigma$ 到 $\sigma$ 的扇的映射，因为 $\Sigma$ 的每个锥都含于 $\sigma$ 中，于是得到态射
\[
  \pi:\tilde X:=X_\Sigma\to X_\sigma=X.
\]
作为环面几何中的练习，可以验证该映射紧合，且诱导同构 $\tilde X\setminus D_\rho\cong X\setminus\{(0,0)\}$。该映射称为环面爆破或加权爆破；若 $m'=(1,1)$，它恰是 $\mathbb{A}^2_k$ 在原点处的通常爆破。（见例 3.3。）

令 $\tilde D\subseteq\tilde X$ 是 $\tilde X$ 上环面除子之并；注意 $\tilde D=\pi^{-1}(D)$。于是得到 $\tilde X$ 上的对数结构 $\cM_{(\tilde X,\tilde D)}$，以及对数态射 $\pi:\tilde X^\dagger\to X^\dagger$，因为 $X$ 上只在 $D$ 上消没的函数拉回到 $\tilde X$ 上就是只在 $\tilde D$ 上消没的函数。

% [原文 p.104]
现在给定如上所述的 $f:\operatorname{Spec} k^\dagger\to X^\dagger$，我们问：该映射何时能提升为交换图表
\[
  \begin{array}{ccc}
    & & \tilde X^{\dagger}\\[6pt]
    \tilde f\nearrow & & \pi\big\downarrow\\[6pt]
    \operatorname{Spec}k^{\dagger} & \xrightarrow{f} & X^{\dagger}
  \end{array}
\]
使得 $\operatorname{Spec} k^\dagger$ 的像落在 $D_\rho\setminus\operatorname{Sing}(\tilde D)$ 中？（这是 $D_\rho$ 的开环面轨道。）

为回答这个问题，先考虑给出映射 $\tilde f$ 意味着什么。这与给出映射 $f$ 的方式相同。选取 $N$ 的一组基 $n_1,n_2$，使得 $\langle n_1,m'\rangle=1$ 且 $\langle n_2,m'\rangle=0$，于是锥 $\rho^{\vee}$ 由 $n_1$ 与 $\pm n_2$ 生成，且 $z^{n_1}$ 沿 $D_\rho$ 消没到一阶。我们指定 $\tilde f$ 的像为点 $x\in D_\rho\setminus\operatorname{Sing}(\tilde D)$，此时 $\cM_{(\tilde X,\tilde D),\bar x}$ 在局部由形如 $\varphi\cdot z^{an_1}$ 的函数组成，其中 $a\in\N$，$\varphi$ 在 $\bar x$ 的邻域内可逆。于是映射 $\tilde f^{\#}$ 由
\[
  \tilde f^{\#}(\varphi\cdot z^{an_1})=(\varphi(x)\tilde g(a),ab)\in k^{\times}\oplus\N
\]
给出，其中 $\tilde g:\N\to k^{\times}$，$b\in\N$，$b\ne 0$。

为使 $\pi\circ\tilde f$ 与 $f$ 一致，需要 $\tilde f^{\#}\circ\pi^{\#}=f^{\#}$。现在对
\[
  \varphi\cdot z^{\beta_1n_1+\beta_2n_2}\in\cM_{(X,D),\bar 0},
\]
我们有
\[
  \begin{aligned}
    \tilde f^{\#}\circ\pi^{\#}(\varphi\cdot z^{\beta_1n_1+\beta_2n_2})
    &=\tilde f^{\#}\bigl((\varphi\circ\pi)\cdot z^{\beta_1n_1+\beta_2n_2}\bigr)\\
    &=\bigl(\varphi(\pi(x))\cdot z^{\beta_2n_2}(x)\cdot\tilde g(\beta_1),\beta_1b\bigr).
  \end{aligned}
\]
这里 $z^{\beta_2n_2}(x)$ 表示有理函数 $z^{\beta_2n_2}$ 在 $x$ 处的值。注意 $z^{n_2}$ 实际上是 $D_\rho$ 上的坐标。由前面 $f^{\#}$ 的公式，我们首先看出需要对所有满足 $\beta_1n_1+\beta_2n_2\in\sigma^{\vee}$ 的 $\beta_1,\beta_2$ 都有
\[
  \beta_1b=\beta_1\langle n_1,m\rangle+\beta_2\langle n_2,m\rangle.
\]
所以特别地必须有 $\langle n_2,m\rangle=0$，即 $m$ 与 $m'$ 成比例，且 $b=\langle n_1,m\rangle$。故必须有 $m=bm'$。其次，需要
\[
  g(\beta_1n_1+\beta_2n_2)=(z^{n_2}(x))^{\beta_2}\cdot\tilde g(\beta_1).
\]
于是特别地，$g$ 决定了 $z^{n_2}(x)$，从而决定了点 $x$ 本身。

所以从哲学上说，对数点的选择指定了一个环面爆破，连同该环面爆破的例外除子上的一个点。还有一些剩余信息，即 $b$ 与 $\tilde g(1)$ 的选择。不过在例 3.17 中，我们将考虑一个使这一额外自由度消失的情形。
\end{example}

% [Tutor Note: p.104 提升判据的核对锚点：$n_1,n_2$ 的选取（$\langle n_1,m'\rangle=1,\langle n_2,m'\rangle=0$）使 $\rho^\vee$ 由 $n_1,\pm n_2$ 生成、$z^{n_1}$ 沿 $D_\rho$ 一阶消没，从而 $D_\rho$ 上可用 $z^{n_2}$ 作坐标；由 $\tilde f^{\#}\circ\pi^{\#}=f^{\#}$ 比较 $\beta$-系数即得 $m=bm'$ 与 $\beta_1b=\langle\beta_1n_1+\beta_2n_2,m\rangle$。]

\begin{example}
其次考虑从对数概形 $X^\dagger$ 到标准对数点的态射 $X^\dagger\to\operatorname{Spec} k^\dagger$。首先 $X$ 是 $k$ 上的概形，且有映射 $f^{\#}:k^{\times}\times\N\to\cM_X$。映射 $k^{\times}\to\cM_X$ 完全确定，由 $c\mapsto\alpha_X^{-1}(c)$ 给出。因此 $f^{\#}$ 完全由映射 $\N\to\cM_X$ 确定。后者又由给出截面 $\rho\in\Gamma(X,\cM_X)$（它是 $1\in\N$ 的像）确定。注意我们要求 $\alpha_X(\rho)=f^*(\alpha_{\operatorname{Spec}k^\dagger}(0,1))=0$。

所以标准对数点上的对数概形，就是 $k$ 上的对数概形连同一个满足 $\alpha_X(\rho)=0$ 的 $\cM_X$ 的截面 $\rho$。
\end{example}

\begin{example}
回到例 3.15 的情形，考虑由单项式 $z^{(1,1)}$ 给出的映射 $\pi:X\to\mathbb{A}^1_k$。$0\in\mathbb{A}^1_k$ 上的纤维恰是 $D$。我们有除子对数结构 $\cM_{(\mathbb{A}^1_k,0)}$，以及对数态射 $\pi:X^\dagger\to(\mathbb{A}^1_k)^\dagger$，这仍是因为 $\mathbb{A}^1_k$ 上在 $0$ 之外可逆的函数拉回到 $X$ 上就是 $D$ 之外可逆的函数。

